% Generated by sync_speaker_notes.py from the presentation outline.
% Edit the outline then rerun the script to refresh these notes.
\newcommand{\SlideOneNotes}{%
\textbf{Diagnostic selection and shared resources}\par\smallskip
\textbf{Time: 0:00--0:45}\par\medskip
Helion wants the on-duty engineer to choose the next diagnostic procedure while accounting for cost and useful completeness. Our research now tests three links: predicting faults, selecting procedures and scheduling shared resources. Better prediction alone need not change the work required. A real coexisting fault remains useful even when it adds work. The engineer retains closure authority, with a proposed coordinator role for dispatch. The headline finding is that the model predicts some synthetic faults usefully, but adds no demonstrated diagnostic value over CT-first rules in this experiment. Scheduling changes milestone feasibility in the bounded example.\par\medskip
{\fontsize{7}{8.5}\selectfont\color{Muted}Evidence: Client brief; design notebook, Helion context; benchmark report.\par}
}
\newcommand{\SlideTwoNotes}{%
\textbf{Implemented pipeline and fair comparators}\par\smallskip
\textbf{Time: 0:45--1:45}\par\medskip
Acceptance has already rejected the stack. Diagnostic findings inform later process investigation without proving which machine caused a defect. We implemented six policies with the same costs, report scope and conservative closure standard. The mock performs triggered branches first. CT-first can advance completeness work. The heuristic ranks eligible tests by inconclusive-adjusted predicted unresolved coverage per resource dollar. It is not an optimal continuation planner. Findings change evidence and eligibility while initial probabilities remain fixed. A low probability never confirms absence. A gross-delamination CT negative cannot exclude all delamination, and IR provides localisation only. Open intact-sample obligations or contradictions block destructive SEM. These restrictions let us compare policies without rewarding omitted necessary work.\par\medskip
{\fontsize{7}{8.5}\selectfont\color{Muted}Evidence: Pipeline README; MOCK-ENG-001; benchmark report, section 2.\par}
}
\newcommand{\SlideThreeNotes}{%
\textbf{Evidence and operating assumptions}\par\smallskip
\textbf{Time: 1:45--3:00}\par\medskip
We retained the 653, 126 and 137 rejected-stack partitions with lot/time separation. All training targets come from supplied synthetic truth. Preprocessing fits training data only. IDs, later annotations, simulator internals and timing margin are excluded. Even so, permitted acceptance readings contain shortcuts: the detected-interconnect count equals the generated TSV-plus-microbump fault count. The manufacturing-only comparison exposes dependence on those inputs. Test data was already inspected, so findings are retrospective. Six crack cases and two coexisting cases limit conclusions. Staff, dollar rates and report-error rates are assumptions. Costs charge technician and engineer labour, equipment occupancy and supplies. They represent consumed resources, not wholly avoidable cash. Forty cases received an independent synthetic audit assignment, fixed across policies. Real independent complete audits remain missing.\par\medskip
{\fontsize{7}{8.5}\selectfont\color{Muted}Evidence: validation.json; SYN-OPS-001; benchmark evidence boundary.\par}
}
\newcommand{\SlideFourNotes}{%
\textbf{Fault-prediction findings}\par\smallskip
\textbf{Time: 3:00--4:30}\par\medskip
We fitted regularised logistic regression with training-only imputation, missingness indicators, scaling and categorical encoding. All learned alternatives used the same training-lot folds and three-value regularisation search. A single saved artifact retains the research alternatives. Each has seven binary heads, so faults may coexist. Average precision summarises precision and recall across ranking thresholds. Macro averaging gives every mechanism equal weight. It differs from precision at a selected threshold and from diagnosis accuracy. Both acceptance-based variants exceed the 0.145 prevalence reference, but full context does not improve the point estimates. The paired 95\% interval for full-minus-inspection macro AP spans $-$0.058 to +0.047, so the small AP difference is uncertain. Full-model Brier score is 0.081 versus inspection's 0.077, with lower preferred. Calibration was assessed without fitting an extra calibrator. Strong electrical prediction and weak rare-fault prediction should not disappear inside one average.\par\medskip
{\fontsize{7}{8.5}\selectfont\color{Muted}Evidence: predictive\_metrics.csv; predictive\_paired\_comparisons.csv; training.json; calibration figure.\par}
}
\newcommand{\SlideFiveNotes}{%
\textbf{Diagnostic selection versus the baselines}\par\smallskip
\textbf{Time: 4:30--6:15}\par\medskip
Policies received the same potential reports for each sample and procedure, with shared audit assignments. A hundred replications vary assumed report outcomes on the same cases. They do not create new independent samples. The full model and CT-first produce identical key outcomes despite different procedure order in 5,796 paired runs. Prevalence and the other model heuristics also match the CT-first aggregate outcomes. Consequently the tiny difference from the mock cannot be credited to ML. Under this catalogue, CT already earns a high coverage-per-dollar score, and changing order often leaves the necessary procedure set unchanged. CT, acoustic, electrical and SEM once each still cost \$2,720. Evidence-complete reports can be wrong, which explains the lower correctly-complete rate. Pending dollars cover known unattempted procedures only. Manual review and future resolution remain unpriced. One-attempt limits and blocked destructive work explain much of the unresolved work. The 1,000 paired lot-cluster bootstrap resamples quantify case uncertainty separately from Monte Carlo variation. Cost, report-error and outage stresses test assumptions, without establishing real-world benefit.\par\medskip
{\fontsize{7}{8.5}\selectfont\color{Muted}Evidence: decision\_metrics.csv (base/test); decision\_uncertainty.csv; verification.json; benchmark report.\par}
}
\newcommand{\SlideSixNotes}{%
\textbf{Bounded scheduling findings}\par\smallskip
\textbf{Time: 6:15--7:30}\par\medskip
The scheduler uses the existing one-day snapshot, with commitments before nine preserved. B needs its warpage question addressed by ten. A's underfill milestone is eleven. Putting B first meets both milestones if CT gives conclusive evidence. TECH-01 prepares each stack for fifteen minutes, and QE-01 interprets the final fifteen minutes. These phases and equipment reservations do not conflict. A-first misses B's milestone, although both orders cost \$240. This demonstrates feasible dispatch against the documented reference, while actual incumbent dispatch remains unknown. An inconclusive B stays unresolved and does not cancel A's accepted booking. In Case B, \$920 is spent after CT, acoustic and electrical work, with SEM pending. Case C's \$2,720 conclusive path remains a conditional cost trace because the supplied calendar cannot support a future SEM appointment.\par\medskip
{\fontsize{7}{8.5}\selectfont\color{Muted}Evidence: scheduling.json; walkthroughs.json; SYN-OPS-001 Case D.\par}
}
\newcommand{\SlideSevenNotes}{%
\textbf{Proposed operation, monitoring and fallback}\par\smallskip
\textbf{Time: 7:30--9:00}\par\medskip
The diagram remains the proposed fab architecture. Our implementation is an offline CLI with saved preprocessing, recorded seeds, dependencies and checksums. It installs no nightly job, interface or production service. Proposed nightly scoring fits the local environment, while findings and resource events would update eligibility and uncommitted bookings between imports. Failed imports produce no fresh scores. Missing required records or unknown tool/supplier conditions route to rules or manual review, and stale calendars prevent new slot promises. Monitoring would separate input drift from audited outcome deterioration. Queue age, inconclusive frequency and engineer overrides are early proxies while labels arrive. Yield Engineering would own outcome review, lab operations the calendars, and IT ingestion alerts. New tools or suppliers trigger qualification review. Independent audits stay protected. Only qualified, complete audit findings enter the reference-label cohort. Inconclusive or contradictory audits remain incomplete. Only reviewed candidates could replace a qualified model. These operational controls still require real records and prospective testing.\par\medskip
{\fontsize{7}{8.5}\selectfont\color{Muted}Evidence: Design notebook, operations; helion-architecture.mmd; pipeline README; run\_manifest.json.\par}
}
\newcommand{\SlideEightNotes}{%
\textbf{Interpretation and next experiments}\par\smallskip
\textbf{Time: 9:00--10:00}\par\medskip
Improving the 0.497 macro AP alone would not establish business benefit. Investigate errors for cracks, voids and delamination, then test limited feature or regularisation changes using training-lot validation. Avoid repeated optimisation against this already-inspected test cohort. In parallel, engineers must identify where qualified alternative procedures could change completion cost. Lab operations and finance should replace invented calendars and rates. Quality must confirm report scope, repeat rules and closure. Roughly forty-five stacks per quarter would be selected for full-battery audits, with fewer potentially yielding complete supported labels. Rare-fault evidence will be limited, so pool history carefully and retain a fresh prospective evaluation. Any subsequent shadow study and pilot need qualification first. The experiment supports honest uncertainty and a continued rules-based option, rather than an automatic model release.\par\medskip
{\fontsize{7}{8.5}\selectfont\color{Muted}Evidence: Benchmark limitations and next evidence; client brief; SYN-OPS-001 section 6.\par}
}
